\subsection{Excision and the norm class with ramification only above p}
\label{sec:kt-cubic-excision-norm}

Retain the hypotheses, coefficients, and specified ordinary quadratic
family of KT-ORDINARY-CUBIC-MASSEY-CONDITION.  In particular \(p>3\),
\(e\) is even, \(\psi=\omega^e\), \(E=\mathbf Q(\sqrt D)\),
\(p\nmid D\), and
\(S=\{\infty,p\}\cup\{q:q\mid D\}\).
Write \(F=\mathbf Q(\mu_p)\), \(K=EF\), and
\(\delta=\omega^{1-e}\).  For a number field \(J\), the notation
\(G_{J,p}\) means its Galois group unramified outside the primes above
\(p\) and the infinite primes.  Let \(A_F\) denote the \(p\)-primary
ideal class group of \(F\).  This section proves an identification of
the actual relative obstruction, not just an equality of dimensions.

\begin{theorem}[KT-CUBIC-EXCISION]
For every finite coefficient field \(k/\mathbf F_p\) there is a
natural isomorphism
\[
 \mathcal O_\psi\ \simeq\ H^2(G_{\mathbf Q,p},k(\psi)).
                                                        \tag{CE.1}
\]
For \(k=\mathbf F_p\), fixing a primitive root \(\zeta\in\mu_p\)
then gives an isomorphism
\[
 \mathcal C_\zeta:\mathcal O_\psi
      \xrightarrow{\ \sim\ }(A_F/pA_F)^{\delta}.
                                                        \tag{CE.2}
\]
There is no quotient in (CE.2) by classes of primes above \(D\).
For larger finite \(k\), (CE.2) extends by scalars.
\end{theorem}

\begin{proof}
Put \(U=\operatorname{Spec}\mathbf Z[1/p]\),
\(V=\operatorname{Spec}\mathbf Z[1/pD]\), and \(M=k(\psi)\).
The corresponding finite locally constant sheaf extends to \(U\).
For \(q\mid D\), let \(U_q=\operatorname{Spec}(\mathbf Z_{(q)})^h\),
with fraction field \(\mathbf Q_q^h\) and residue field \(\mathbf F_q\).
Etale excision supplies the distinguished triangle
\[
 R\Gamma(U,M)\longrightarrow
 R\Gamma(V,M)\oplus\bigoplus_{q\mid D}R\Gamma(\mathbf F_q,M)
 \longrightarrow\bigoplus_{q\mid D}R\Gamma(\mathbf Q_q,M).
                                                        \tag{CE.3}
\]
The last arrow is the difference of restriction and unramified
inflation.  Here is the precise source of this local gluing statement.
For \(j:V\hookrightarrow U\), the complex form of Milne's
Lemma II.2.4 gives
\(R\Gamma(U,j_!M)\to R\Gamma(V,M)\to
\bigoplus R\Gamma(\mathbf Q_q^h,M)\).
The exact sequence \(0\to j_!M\to M\to i_*i^*M\to0\)
gives the second triangle with the same first term.
Comparing them gives (CE.3).
Milne II.1.1(b) identifies cohomology over \(U_q\) with that over
its residue field; the henselian fraction field and its completion
have equivalent finite separable extensions and the same absolute
Galois group, as recalled explicitly in the proof of I.2.14, p.~37.
Thus the final terms are exactly the completed local groups used
in (CM.14), with their actual unramified inflation maps.

One can read the degree-two group of the mapping fibre in (CE.3)
without discarding any local data.  A representative consists of
a global two-cocycle \(c\), unramified local two-cocycles \(b_q\),
and local one-cochains \(z_q\) with
\(dz_q=c|_q-b_q\).
Since \(H^2(\mathbf F_q,M)=0\), subtract primitives of the \(b_q\)
and set them to zero.  The remaining ambiguity in those primitives
is exactly unramified local \(H^1\), in addition to local
coboundaries.  The resulting quotient is therefore the pairs
\((c,(z_q)_{q\mid D})\) of (CM.14), with precisely its stated
equivalence and its allowed spaces \(L_q\).

Forgetting the \(p\)-primitive in (CM.14) is an isomorphism on
that quotient: local duality gives
\(H^2(\mathbf Q_p,M)=H^0(\mathbf Q_p,k(\delta))^\vee=0\),
because \(1-e\) is odd and is not zero modulo \(p-1\), while
\(L_p\) is all local \(H^1\).  The infinite place has no
positive-degree \(p\)-cohomology.  Finally Milne II.2.9 applies
on both affine schemes, since \(p\) is invertible, and identifies
their etale cohomology with the corresponding Galois cohomology.
This proves (CE.1) with its natural map, rather than by comparing
dimensions.

For (CE.2), the cover
\(U_F=\operatorname{Spec}\mathcal O_F[1/p]\to U\) is finite
etale of degree \(p-1\).  Exact averaging gives
\[
 H^2(U,\mathbf F_p(\psi))
     \simeq H^2(U_F,\mathbf F_p)^{\psi^{-1}}.
                                                        \tag{CE.4}
\]
The identification \(1\mapsto\zeta\) takes the right side to
\(H^2(U_F,\mu_p)^\delta\): it contributes the additional
cyclotomic character \(\omega\).  The Kummer exact sequence is
\[
 0\longrightarrow \operatorname{Pic}(U_F)/p
 \longrightarrow H^2(U_F,\mu_p)
 \longrightarrow \operatorname{Br}(U_F)[p]\longrightarrow0.
                                                        \tag{CE.5}
\]
There is one prime of \(F\) above \(p\), and \(F\) has no real
places.  The local invariant sequence for the Brauer group shows
\(\operatorname{Br}(U_F)[p]=0\): a single local invariant must
have sum zero.  The prime above \(p\) is principal, generated
by \(1-\zeta\), so \(\operatorname{Pic}(U_F)=\operatorname{Cl}(F)\).
Consequently (CE.4)--(CE.5) give exactly (CE.2).  Extending the
continuous cochain complexes by the finite field \(k\) proves
the scalar extension statement.
\end{proof}

\begin{lemma}[KT-CUBIC-MATCHED-NORM-SYSTEM]
Let \(J\) be a field of characteristic zero containing \(\mu_p\),
let \(\Gamma\) be a quotient of \(G_J\), and let
\(\chi:\Gamma\to H=\mathbf F_p\) be surjective.
Set \(\mathcal R=\mathbf F_p[H]=\mathbf F_p[x]/(x^p)\),
where \(x=[\tau]-1\) and \(\tau\) corresponds to \(1\in H\),
and write \(I=(x)\).
Suppose \(f_2\in Z^1(\Gamma,\mathcal R/I^2)\), with augmentation
\(v\), has a Bockstein class whose inflation to
\(H^2(G_J,\mathbf F_p)\) is zero.  Let \(L/J\) be cut out
by \(\chi\), and fix a Kummer representative \(B_0\in J^\times\)
for \(v\), using \(\zeta\).

There is a norm representative \(y\in L^\times\),
\(N_{L/J}y=B_0\), for which Sharifi's truncated norm system
represents the specified class \([f_2]\).  In particular its
proper triple value is the actual Bockstein of \(f_2\),
not merely its coset modulo lower products.  If
\(\Gamma=\operatorname{Gal}(\Omega/J)\), then
\(D^{(1)}y\in\Omega^{\times p}\).
\end{lemma}

\begin{proof}
All lifting in this paragraph takes place on the absolute
group \(G_J\).  The cyclic norm-residue sequence is exact at
its middle degree-two term:
\[
 H^1(G_J,\mathbf F_p)
 \xrightarrow{\chi\smile-}H^2(G_J,\mathbf F_p)
 \xrightarrow{\operatorname{res}}H^2(G_L,\mathbf F_p).
                                                        \tag{CE.9}
\]
This is the absolute Galois input stated in
Lam--Liu--Sharifi--Wake--Wang, Definition~6.1.2 and the
following paragraph, p.~48; their Proposition~6.1.5 uses
the same coefficient correction as below.  We do not
apply (CE.9) to a group with restricted ramification.

The assumed zero Bockstein first lifts the inflated
\([f_2]\) to \(H^1(G_J,\mathcal R/I^3)\).
For \(3\le n<p\), suppose \(f_n\) is a cocycle lifting it
through \(I^n\), and let \(\alpha_n\) be the obstruction
to lifting through \(I^{n+1}\).  On \(G_L\) the entire
regular module is trivial, so that coefficient sequence
splits there.  Thus \(\operatorname{res}\alpha_n=0\).
By (CE.9), write \(\alpha_n=\chi\smile\gamma_n\)
for a scalar one-cocycle \(\gamma_n\).
Replacing \(f_n\) by \(f_n-\gamma_n x^{n-1}\)
kills the obstruction and preserves its image modulo
\(I^2\).  The sign is explicit: for the differential used
throughout this report,
\[
 d(\gamma_n x^{n-1})(g,h)
       =\chi(g)\gamma_n(h)x^n\pmod{x^{n+1}}.
                                                       \tag{CE.10}
\]
The scalar coefficient is a cocycle and
\((1+x)^{\chi(g)}=1+\chi(g)x+O(x^2)\), which proves
(CE.10).  Iterating gives
\(F_0\in Z^1(G_J,\mathcal R)\) whose image in
\(H^1(G_J,\mathcal R/I^2)\) is the inflated \([f_2]\).
If the representative differs by a coboundary, lift
its zero-cochain to \(\mathcal R\) and subtract that
coboundary.  Its projection can therefore be taken
literally equal to the inflated \(f_2\).

Shapiro identifies \([F_0]\) with a Kummer class
\(y\in L^\times/L^{\times p}\); augmentation is
corestriction, so \(N_{L/J}y\equiv B_0\pmod{J^{\times p}}\).
We spell out the truncated-system identification.
Evaluate \(F_0|_{G_L}\) at the coefficient of the
identity of \(H\).  This is the Kummer character of
\(y\).  Equivariance makes its coefficient at
\(\sigma^j\) the Kummer character of \(\sigma^j y\).
Thus its image in the regular module is
exactly the map \(\phi(\lambda)=\sum_h\lambda(h)h\)
of Sharifi's Lemma~3.1, with its stated inversion
pairing, used in (4.2) and Lemma~4.1.
Projecting to \(\mathcal R/I^3\) and taking the action
\((1+x)^\chi\) gives the binomial four-by-four system;
its two last-column entries modulo the central entry
are precisely the two coefficients of \(f_2\).

One normalization in Lemma~4.1 causes no discrepancy.
For a lift \(\sigma\in G_J\) with \(\chi(\sigma)=1\),
\(F_0(\sigma)-v(\sigma)\) lies in \(I=x\mathcal R\).
Conjugation by a constant element of the additive
regular module removes that difference, since its
coboundary at \(\sigma\) is multiplication by \(x\).
This puts the image of \(\sigma\) into Sharifi's
\(X+v(\sigma)E_{34}\) form.  It changes \(f_2\)
only by a coboundary and preserves its Bockstein class.
The same conjugation descends on \(\Gamma\), since
the action \(\chi\) does.  This verifies matching of
the specified truncated class, rather than just
existence of some proper system with first entries
\(\chi,\chi,v\).

Choose an actual representative \(y\).  If its norm is
\(B_0c^p\), divide \(y\) by the base-field element \(c\).
Under Shapiro this subtracts a cocycle supported on
the norm element
\(\sum_{h\in H}[h]=x^{p-1}\), so it does not change
the projection modulo \(I^2\).  The norm is now
exactly \(B_0\).  Finally the matched system modulo
\(I^2\) factors through \(\Gamma\), so on \(G_\Omega\)
the full regular cocycle has values in \(I^2\).
By Sharifi's Lemma~3.3, \(D^{(1)}\) belongs to
\((I^2)^\perp\) for precisely the pairing just used.
The Kummer character of \(D^{(1)}y\) consequently
vanishes on \(G_\Omega\).  This proves
\(D^{(1)}y\in\Omega^{\times p}\) and the required
proper norm condition.
The full regular cocycle was constructed only on
\(G_J\), and its coefficients of degree at least two
may ramify outside the set allowed by \(\Gamma\).
Only its specified truncation is required to descend.
In particular \(D^{(2)}y\) need not have a \(p\)-th
root in \(\Omega\); the restricted-ramification
transgression being computed need not vanish.
\end{proof}

\begin{theorem}[KT-CUBIC-P-ONLY-NORM]
Assume \(k=\mathbf F_p\) and \([u]\ne0\).  Restrict the specified
ordinary quadratic family to \(K=E(\mu_p)\).
Its binomial defining system (CS.9), after the stated conjugation,
factors through \(G_{K,p}\).  Thus the proper norm construction
of KT-CUBIC-BOCKSTEIN-NORM-IDENTIFICATION can be made with
\(S_K\) consisting only of the primes above \(p\).

Let \(L/K\) be the degree-\(p\) extension cut out by \(u|_{G_{K,p}}\),
and use this smaller ramification set in its proper norm system.
Let \(\mathfrak Y\) and \(P^{(1)}_{1,L/K,p}\) be Sharifi's
ideal and lower-product subgroup for this system.  There is a
specified projection
\[
 \pi_\delta:H^2(G_{K,p},\mu_p)
           \longrightarrow (A_F/pA_F)^\delta
                                                        \tag{CE.6}
\]
which is inverse to extension on the \(E\)-positive
\(\delta\)-component.  The actual ordinary obstruction satisfies
\[
 \mathcal C_\zeta(\mathfrak o_3(u,v))
 \equiv -3\pi_\delta([N_{L/K}\mathfrak Y])
          \pmod{\pi_\delta(P^{(1)}_{1,L/K,p})}.
                                                        \tag{CE.7}
\]
In particular a nonzero right-hand coset proves
\(\mathfrak o_3(u,v)\ne0\).  All local correction information
in (CM.15) is retained by (CE.2) and this smaller ramification
set.  Moreover one may choose \(y\) to match the specified
ordinary system by KT-CUBIC-MATCHED-NORM-SYSTEM.  For that
matched choice, the ideal furnished by Sharifi's proof satisfies
the stronger equality
\[
 \mathcal C_\zeta(\mathfrak o_3(u,v))
            =-3\pi_\delta([N_{L/K}\mathfrak Y])
       \quad\hbox{in }(A_F/pA_F)^\delta .
                                                       \tag{CE.11}
\]
The quotient in (CE.7) is necessary when a proper norm
representative is used without retaining that matching.
\end{theorem}

\begin{proof}
At a prime \(q\mid D\), the actual ordinary quadratic family
is unramified after restriction to \(E\).  This was proved in
KT-GLOBAL-QUADRATIC-MIXED-PERSISTENCE: locally it is conjugate to
\(\exp(t v(g)Y)r_0(g)\), and both terms are one on inertia
over \(E\).  Identity on inertia is preserved by that fixed
conjugation.  Over \(K\), all residual characters are trivial.
Hence the restricted quadratic homomorphism is one on every
inertia group away from \(p\), and factors through \(G_{K,p}\).
Its logarithmic coordinates \(u,v,w\), and therefore the
polynomial binomial entries in (CS.9), factor through that
same quotient.  This statement uses the actual local family,
not just the vanishing of its inertia cohomology classes.

First restrict to \(G_{K,S_K^{\rm all}}\), where
\(S_K^{\rm all}\) contains the primes above \(pD\).
This is the finite-index subgroup of \(G_{\mathbf Q,S}\),
so restriction is injective in degree one by restriction
and corestriction, because \([K:\mathbf Q]\) is prime to \(p\).
Then descend through its quotient \(G_{K,p}\).
The descended \(\mathbf F_p\)-valued
character \(u\) remains nonzero and therefore cuts out
a degree-\(p\) extension unramified outside \(p\).
The binomial Bockstein calculation (CS.11)--(CS.14)
therefore applies on \(G_{K,p}\) itself.  The existence
of a matched norm representative will follow below;
it is not inferred solely from the terminology
``proper defining system.''

We next check that this is the image of the \emph{relative}
class in (CE.1).  The excision construction (CE.3) is functorial
for the finite map
\(U_K=\operatorname{Spec}\mathcal O_K[1/p]\to U\),
even though that map is ramified at primes over \(D\).
Use the corresponding local gluing triangles before and after
this base change.  Its pullback on a relative representative
restricts the global cocycle and all the local primitives.
At a prime above \(D\), the chosen acceptable cubic local
lift is unramified over \(K\).  Thus its central degree-three
coordinate is an unramified primitive of the pulled-back
\(-\Omega_+\).  Since the quadratic coordinates already
factor through \(G_{K,p}\), this is precisely the relative
representative of their raw cubic obstruction on \(G_{K,p}\).
Changing either primitive makes an unramified local cocycle,
which is killed by the equivalence in (CE.3).
There is no lost primitive at \(p\): it was forgotten on the
\(\mathbf Q\)-side because \(L_p\) was full.  On the
\(K\)-side this paragraph uses the gluing complex only at
primes over \(D\), so it does not incorrectly impose a
vanishing condition on \(H^2(K_w,\mathbf F_p)\) at \(w\mid p\).
It follows that the image of (CE.1) on \(U_K\) is exactly
the raw cubic class of the descended defining system.

For completeness that pullback does not lose the rational
class in (CE.2).  The fields \(E\) and \(F\) are disjoint,
since \(p\nmid D\), and
\(\operatorname{Gal}(K/\mathbf Q)
  =\operatorname{Gal}(E/\mathbf Q)\times\operatorname{Gal}(F/\mathbf Q)\).
Let \(e_{+,\delta}\) be the averaging projector onto the
component trivial on the first factor and equal to \(\delta\)
on the second.  It is defined over \(\mathbf F_p\), since
the group order \(2(p-1)\) is a unit.
The classes of primes of \(K\) above \(p\), and the local
Brauer invariant module at those primes, are permutation
modules on which the cyclotomic factor acts trivially.
Indeed \(F/\mathbf Q\) is totally ramified at \(p\), while
\(E/\mathbf Q\) is unramified there.  Since \(\delta\ne1\),
neither module has a \((+,\delta)\)-component.  Kummer and
the Brauer local invariant sequence therefore identify
\[
 e_{+,\delta} H^2(G_{K,p},\mu_p)
       \simeq (\operatorname{Cl}(K)/p)^{{+,\delta}}.
                                                        \tag{CE.8}
\]
Extension and norm on ideal classes have compositions
\(N_{K/F}\circ\mathrm{ext}=2\) and
\(\mathrm{ext}\circ N_{K/F}=1+c\), where
\(c\) generates \(\operatorname{Gal}(K/F)\).
They identify the right side of (CE.8) with
\((A_F/pA_F)^\delta\).  Define \(\pi_\delta\) by first
applying \(e_{+,\delta}\), then (CE.8), then
\(\tfrac12N_{K/F}\).  This proves (CE.6), including
injectivity of the pullback on the image of (CE.2).
There is no need to infer this injectivity from a putative
unramified Galois cover at the primes over \(D\).

The raw ordinary class on \(G_{K,p}\) lies in the
Picard subgroup by (CE.8).  Its inflation to the
absolute group \(G_K\) is therefore zero: the
restriction of a line bundle to \(\operatorname{Spec}K\)
is trivial, and the Kummer boundary is natural.
By (CS.10)--(CS.11), its value is three times the
Bockstein of the specified \(f_2=v-ax\).  Since
\(3\) is invertible, that Bockstein also inflates
to zero.  KT-CUBIC-MATCHED-NORM-SYSTEM now provides
a genuinely matched \(y\) over the absolute group,
whose truncated system descends to \(G_{K,p}\).

For this particular \(y\), Sharifi's Proposition~4.2
identifies the value of its specified defining
system exactly with \(\operatorname{Tra}[D^{(2)}y]\).
The proof of Theorem~5.4, equation (5.2), identifies
this exact transgression with \(-[N_{L/K}\mathfrak Y]\).
The possible common prime-to-\(p\) power taken in
that proof can be canceled in the \(\mathbf F_p\)-coefficient
group.  There is no boundary term because
\(2<p-1\).  The lower-product quotient is introduced
in the statement of the theorem to forget which
proper system was chosen; it is unnecessary for
these two identities attached to the same \(y\).
The exact cochain identities (CS.3)--(CS.10) say
that the raw ordinary obstruction is three times
this binomial-system value.  Applying \(\pi_\delta\)
therefore proves (CE.11).  If an arbitrary proper
norm representative is used, Sharifi's Theorem~5.4
instead gives its coset modulo lower products,
and proves (CE.7).
The notation \(\pi_\delta(P)\) denotes the image subgroup;
no stability assertion about \(P\) is needed for this
projection.  A nonzero resulting coset is a nonzero full
class, so (CE.2) proves the final implication.
\end{proof}

\paragraph{What this comparison changes, and what remains.}
The raw formula with \(S_K\) containing primes above \(D\)
could erase an ordinary obstruction supported on those
ideal classes.  Equations (CE.1)--(CE.11) remove that loss:
the specified obstruction is an element of the full
rational \(\omega^{1-e}\) class component, and the
proper norm system can be formed with ramification only
above \(p\).  A matched norm representative gives
the exact class (CE.11), but no such norm representative
or ideal class has been explicitly calculated here.
No nonzero value of that element or of its norm coset
has been established.  In particular this
comparison supplies neither an injection from the
original strict \(\omega^{-e}\) Selmer group nor a
source-to-independent-jet theorem.  The larger finite
coefficient-field norm calculation still requires its
own scalar or cyclic reduction.  An arithmetic exclusion
would require an evaluated nonzero norm coset (or a
matched evaluation before the lower-product quotient)
and the separate source-to-jet hypotheses.

\paragraph{Primary inputs and check scope.}
The etale statements used here are Milne,
\emph{Arithmetic Duality Theorems}, second edition (2006):
I.2.14 and its proof, pp.~36--37; II.0.8 and II.0.10,
pp.~145--148; the complete II.1.1, p.~149; II.2.1 and
Remark II.2.2(a), pp.~163--165; II.2.3--2.4 and the complex-level proof,
pp.~166--168; and the complete comparison proposition
II.2.9, pp.~170--172.  These passages supply the local
gluing maps, Kummer and Brauer computations, and the
comparison with the actual finite-ramification Galois
groups.  Milne's official addendum, revised 26 October 2025,
corrects the order of \(i^*\) and \(i_*\) at printed p.~168,
line~5; the sequence above uses the corrected \(i_*i^*\).
The correction is on p.~1 of
\url{https://www.jmilne.org/math/Books/add/ADT2006.pdf}.
The norm matching additionally uses the absolute cyclic
norm-residue exact sequence of Lam--Liu--Sharifi--Wake--Wang,
Definition~6.1.2 and its following paragraph, and the
coefficient correction in the complete Proposition~6.1.5,
pp.~48--49.  The exact Kummer-dual and norm-system
identifications are Sharifi's complete Lemma~3.1,
pp.~13--14, Lemma~4.1, pp.~19--21, Proposition~4.2,
p.~21, and equation (5.2), p.~27.
The proper norm theorem is Sharifi,
\emph{Massey products and ideal class groups},
Theorem~5.4, pp.~26--27 of the 26 July 2006 author
version, with the binomial matrix conventions of
Lam--Liu--Sharifi--Wake--Wang, Section~4.3.  Their
terminology alone is not used to infer a matched norm
representative: that requires the zero absolute
obstruction and KT-CUBIC-MATCHED-NORM-SYSTEM above.
The foundational theorems are imported.  This new proof
checks their scope and the relative comparison; it does
not reconstruct etale excision or the general norm
theorem, and it does not assert an arithmetic
nonvanishing computation.
